% GENERATED FILE. Edit canonical lessons, then run npm run generate:books.
% canonical-source-hash: fnv1a64:b89c10994b5eb763
% canonical-lessons: FR-C459-ecureuil, FR-C459-herisson, FR-C459-sanglier, FR-C459-cerf, FR-C459-biche

\chapter{Les animaux des bois}
\label{ch:fr-les-animaux-des-bois}
\hlchaptermodality{459}

\begin{chapteropening}
\textbf{By the end of this chapter:} \emph{I can name animals of the woods: the squirrel, the hedgehog, the wild boar, the stag and the doe.}

Complete the last lesson of chapter 459: I can name animals of the woods: the squirrel, the hedgehog, the wild boar, the stag and the doe.
\end{chapteropening}

\section[l'écureuil]{l'écureuil — the squirrel}

\label{lesson:FR-C459-ecureuil}



\begin{quote}
Before the new one: say the French for the reef, then the French for the seashell.
\end{quote}

\subsection*{You'll want to know: l'écureuil}
\textbf{l'écureuil} — \textquotedblleft{}the squirrel\textquotedblright{}.

\begin{quote}
Un écureuil monte dans l'arbre avec une noisette.
\end{quote}

\begin{quote}
\emph{A squirrel runs up the tree with a hazelnut.}
\end{quote}

A masculine noun. The ending sounds like \emph{euil}.

\begin{cousinweb}
\emph{écureuil} comes from Middle French \emph{escureul} (EtymDB), from Vulgar Latin \emph{scūriolus}, a small Latin \emph{sciūrus}, from Greek \emph{skíouros} 'shadow-tail'.
\end{cousinweb}

\begin{grammarlens}[title={What you've built}]
One more word for telling, explaining and describing at B1.
\end{grammarlens}

\subsection*{Your turn}
\begin{itemize}
\raggedright
  \item \emph{Say it:} \emph{l'écureuil}
  \item \emph{Say it:} \emph{l'écureuil}, once more
  \item \emph{Say it:} a sentence of your own with \emph{l'écureuil}
  \item \emph{Recall:} say the French for the reef, then the French for the seashell, then say \emph{l'écureuil} again
\end{itemize}

\begin{tcolorbox}[breakable,colback=teal!4,colframe=teal!35!black,title={Before you move on}]
What does l'écureuil mean? (\textquotedblleft{}The squirrel\textquotedblright{}.) Say it once more.
\end{tcolorbox}
\section[le hérisson]{le hérisson — the hedgehog}

\label{lesson:FR-C459-herisson}



\begin{quote}
Before the new one: say the French for the seashell, then the French for the squirrel.
\end{quote}

\subsection*{You'll want to know: le hérisson}
\textbf{le hérisson} — \textquotedblleft{}the hedgehog\textquotedblright{}.

\begin{quote}
La nuit, un hérisson traverse le jardin.
\end{quote}

\begin{quote}
\emph{At night a hedgehog crosses the garden.}
\end{quote}

The \emph{h} is aspirate: \emph{le hérisson}, not \emph{l'hérisson}.

\begin{cousinweb}
\emph{hérisson} comes from Old French \emph{heriçun} (EtymDB), which goes back to Latin \emph{ēricius} 'hedgehog'.
\end{cousinweb}

\begin{grammarlens}[title={What you've built}]
One more word for telling, explaining and describing at B1.
\end{grammarlens}

\subsection*{Your turn}
\begin{itemize}
\raggedright
  \item \emph{Say it:} \emph{le hérisson}
  \item \emph{Say it:} \emph{le hérisson}, once more
  \item \emph{Say it:} a sentence of your own with \emph{le hérisson}
  \item \emph{Recall:} say the French for the seashell, then the French for the squirrel, then say \emph{le hérisson} again
\end{itemize}

\begin{tcolorbox}[breakable,colback=teal!4,colframe=teal!35!black,title={Before you move on}]
What does le hérisson mean? (\textquotedblleft{}The hedgehog\textquotedblright{}.) Say it once more.
\end{tcolorbox}
\section[le sanglier]{le sanglier — the wild boar}

\label{lesson:FR-C459-sanglier}



\begin{quote}
Before the new one: say the French for the squirrel, then the French for the hedgehog.
\end{quote}

\subsection*{You'll want to know: le sanglier}
\textbf{le sanglier} — \textquotedblleft{}the wild boar\textquotedblright{}.

\begin{quote}
Les sangliers sortent du bois le soir.
\end{quote}

\begin{quote}
\emph{The wild boar come out of the wood in the evening.}
\end{quote}

\emph{La laie} is the female.

\begin{cousinweb}
\emph{sanglier} goes back to Latin \emph{singulāris} 'solitary', as in \emph{singulāris porcus} 'the pig that lives alone', from \emph{singulus} 'single' (Lewis \& Short).
\end{cousinweb}

\begin{grammarlens}[title={What you've built}]
One more word for telling, explaining and describing at B1.
\end{grammarlens}

\subsection*{Your turn}
\begin{itemize}
\raggedright
  \item \emph{Say it:} \emph{le sanglier}
  \item \emph{Say it:} \emph{le sanglier}, once more
  \item \emph{Say it:} a sentence of your own with \emph{le sanglier}
  \item \emph{Recall:} say the French for the squirrel, then the French for the hedgehog, then say \emph{le sanglier} again
\end{itemize}

\begin{tcolorbox}[breakable,colback=teal!4,colframe=teal!35!black,title={Before you move on}]
What does le sanglier mean? (\textquotedblleft{}The wild boar\textquotedblright{}.) Say it once more.
\end{tcolorbox}
\section[le cerf]{le cerf — the stag, the deer}

\label{lesson:FR-C459-cerf}



\begin{quote}
Before the new one: say the French for the hedgehog, then the French for the wild boar.
\end{quote}

\subsection*{You'll want to know: le cerf}
\textbf{le cerf} — \textquotedblleft{}the stag, the deer\textquotedblright{}.

\begin{quote}
Au lever du jour, nous avons vu un cerf dans la clairière.
\end{quote}

\begin{quote}
\emph{At daybreak we saw a stag in the clearing.}
\end{quote}

The \emph{f} is silent: \emph{le cerf} rhymes with \emph{serre}.

\begin{cousinweb}
\emph{cerf} is inherited through Old French from Latin \emph{cervus} 'deer' (EtymDB).
\end{cousinweb}

\begin{grammarlens}[title={What you've built}]
One more word for telling, explaining and describing at B1.
\end{grammarlens}

\subsection*{Your turn}
\begin{itemize}
\raggedright
  \item \emph{Say it:} \emph{le cerf}
  \item \emph{Say it:} \emph{le cerf}, once more
  \item \emph{Say it:} a sentence of your own with \emph{le cerf}
  \item \emph{Recall:} say the French for the hedgehog, then the French for the wild boar, then say \emph{le cerf} again
\end{itemize}

\begin{tcolorbox}[breakable,colback=teal!4,colframe=teal!35!black,title={Before you move on}]
What does le cerf mean? (\textquotedblleft{}The stag, the deer\textquotedblright{}.) Say it once more.
\end{tcolorbox}
\section[la biche]{la biche — the doe, the hind}

\label{lesson:FR-C459-biche}



\begin{quote}
Before the new one: say the French for the wild boar, then the French for the stag.
\end{quote}

\subsection*{You'll want to know: la biche}
\textbf{la biche} — \textquotedblleft{}the doe, the hind\textquotedblright{}.

\begin{quote}
La biche et son petit boivent à l'étang.
\end{quote}

\begin{quote}
\emph{The doe and her fawn are drinking at the pond.}
\end{quote}

Affectionate: \emph{ma biche} — my darling.

\begin{cousinweb}
\emph{biche} comes from Middle French \emph{biche} (EtymDB), which is usually traced through a northern dialect to Latin \emph{bestia} 'animal'; the details are not settled.
\end{cousinweb}

\begin{grammarlens}[title={What you've built}]
One more word for telling, explaining and describing at B1.
\end{grammarlens}

\subsection*{Your turn}
\begin{itemize}
\raggedright
  \item \emph{Say it:} \emph{la biche}
  \item \emph{Say it:} \emph{la biche}, once more
  \item \emph{Say it:} a sentence of your own with \emph{la biche}
  \item \emph{Recall:} say the French for the wild boar, then the French for the stag, then say \emph{la biche} again
\end{itemize}

\begin{tcolorbox}[breakable,colback=teal!4,colframe=teal!35!black,title={Before you move on}]
What does la biche mean? (\textquotedblleft{}The doe, the hind\textquotedblright{}.) Say it once more.
\end{tcolorbox}
